\documentclass[11pt]{article}
\usepackage[T1]{fontenc}
\usepackage{lmodern}
\usepackage[margin=1in]{geometry}
\usepackage{amsmath,amssymb,amsthm,mathtools}
\usepackage{microtype}
\usepackage[hidelinks]{hyperref}
\usepackage{enumitem}
\newtheorem{theorem}{Theorem}[section]
\newtheorem{lemma}[theorem]{Lemma}
\newtheorem{proposition}[theorem]{Proposition}
\newtheorem{corollary}[theorem]{Corollary}
\theoremstyle{remark}\newtheorem{remark}[theorem]{Remark}
\newcommand{\Rea}{\operatorname{Re}}
\newcommand{\Ima}{\operatorname{Im}}
\newcommand{\dd}{\mathrm{d}}
\newcommand{\calH}{\mathcal H}
\title{A logarithmic-derivative comparison for Dirichlet Newman families}
\author{Drafted for Edward Baker}
\date{17 October 2026}
\begin{document}
\maketitle
\begin{center}\small Research draft. Proofs and scalar arithmetic have internal same-model LLM checks;\par
independent specialist review and a priority audit remain outstanding.\end{center}
\begin{abstract}
We give an explicit lower bound for the off-zero logarithmic derivative of
a primitive Dirichlet heat deformation. A scalar comparison then controls
the time at which all its zeros become real. A real sine tail added to
canonical-product truncations preserves the required initial lower bound
on the whole real line and allows the comparison to pass to the original
function without a compact-confinement hypothesis. In the normalization
$s=(1+iz)/2$, the argument gives $\Lambda_\chi\le27/100$ for every primitive
character of conductor $q\ge10^{12}$, and
$\Lambda_\chi=O((\log q)^{-4/7})$ uniformly as $q\to\infty$.
The finite scalar inequalities use exact rational enclosures. Extended
proofs and the certificate specification are retained in the accompanying
research notes.
\end{abstract}

\section{Normalization and statements}
Let $\chi$ be primitive of conductor $q>1$, with
$\chi(-1)=(-1)^{\mathfrak a}$, $\mathfrak a\in\{0,1\}$.
Choose a unit phase $\alpha_\chi$ with
$\alpha_\chi^2\varepsilon(\chi)=1$, where
$\varepsilon(\chi)=\tau(\chi)/(i^{\mathfrak a}\sqrt q)$, and set
\begin{align}
 \xi_\chi(s)&=\alpha_\chi(q/\pi)^{s/2}
 \Gamma\bigl((s+\mathfrak a)/2\bigr)L(s,\chi),\label{eq:xi}\\
 \Phi_\chi(u)&=4\alpha_\chi(\pi/q)^{\mathfrak a/2}
 e^{(1+2\mathfrak a)u}\sum_{n\ge1}\chi(n)n^{\mathfrak a}
 e^{-\pi n^2e^{4u}/q},\label{eq:phi}\\
 H_{t,\chi}(z)&=\int_{\mathbb R}e^{tu^2}\Phi_\chi(u)e^{izu}\,\dd u.
 \label{eq:heat}
\end{align}
Then $H_{0,\chi}(z)=\xi_\chi((1+iz)/2)$ and
$\partial_tH_{t,\chi}=-\partial_z^2H_{t,\chi}$.
The Hermitian functional equation makes $H_{t,\chi}$ real entire, although
it need not be even. With compatible phases,
$H_{t,\bar\chi}(z)=\overline{H_{t,\chi}(-\bar z)}$.
The normalization is detailed in Note~1 \cite{note1} and agrees with
Dobner's framework \cite{dobner}. Its associated constant $\Lambda_\chi$
is characterized by reality of all zeros exactly for $t\ge\Lambda_\chi$.

\begin{theorem}[Dirichlet consequence]\label{thm:dirichlet}
Every primitive Dirichlet character of conductor $q\ge10^{12}$ satisfies
\[\Lambda_\chi\le\frac{27}{100}.\]
Moreover, uniformly over primitive characters as $q\to\infty$,
\[0\le\Lambda_\chi=O\bigl((\log q)^{-4/7}\bigr),\qquad
\Lambda_\chi\longrightarrow0.\]
These statements use neither GRH, a quasi-RH half-plane, nor an assumed
confinement theorem for the character flow.
\end{theorem}

The lower bound is Dobner's published theorem \cite{dobner}. The upper
bound will follow from a general comparison. For $C>0$ define
\begin{align}
 D_C(t)&=1+2C(1-e^{-t/2}),&
 m_C(t)&=\frac{Ce^{-t/2}}{D_C(t)},\label{eq:dm}\\
 I_C(t)&=e^{-t}D_C(t)^{4/3},&
 J_C(t)&=\int_0^t I_C(s)\,\dd s.\label{eq:ij}
\end{align}

\begin{theorem}[Comparison without confinement]\label{thm:general}
Let $F$ be a nonzero real entire function of order at most one, with all
zeros in $|\Ima z|\le1$. Suppose
\[ -\Ima\frac{F'}F(x+3i)\ge C>0\qquad(x\in\mathbb R). \]
If $0<T\le1/2$ and $J_C(T)\ge1/2$, then
$e^{-T\partial_z^2}F$ has only real zeros.
\end{theorem}

The proof first treats flows with confined nonreal zeros, then removes
that restriction by a sine-tail approximation. No evenness assumption is
present. The extended approximation argument is in Note~3 \cite{note3}.

\section{Canonical sums and a positive-time lower bound}
For a real entire function with its zeros $\rho=a+iv$ in a bounded strip,
the genus-one product gives, at $p$ above that strip,
\begin{equation}
 S(x,p,t):=-\Ima\frac{F_t'}{F_t}(x+ip)
 =\sum_\rho\frac{p-v}{(x-a)^2+(p-v)^2}.\label{eq:poisson}
\end{equation}
Zeros are counted with multiplicity, including a possible zero at the
origin. The imaginary compensators cancel under conjugation. Their
absolute summability follows from bounded imaginary parts and
$\sum_{\rho\ne0}|\rho|^{-2}<\infty$. A separate, possibly divergent,
real sum of reciprocal zeros is never used.

For the Dirichlet flow, the double-exponential kernel decay gives
\begin{equation}
 \log\max_{|z|\le R}|H_{t,\chi}(z)|
 =O_{\chi,T}((1+R)\log(2+R))\quad(0\le t\le T).
 \label{eq:growth}
\end{equation}
The approximants in Section~\ref{sec:approx} have exponential-type growth.
These bounds also justify the canonical differentiations and completeness
of the real characteristics below. Details, including the zero-free
anchor and uniform Jensen estimate, appear in Note~2, Section~2
\cite{note2}.

\begin{proposition}[Burgers lower bound]\label{prop:floor}
For the Dirichlet flow or the exponential-type approximants below,
suppose the zeros remain in $|\Ima z|\le1$ and $S(x,p,0)\ge C>0$.
Put $d=p-1>0$. Then
\begin{equation}
 S(x,p,t)\ge
 \frac{Ce^{-2t/d^2}}{1+dC(1-e^{-2t/d^2})}
 \qquad(x\in\mathbb R,\ 0\le t\le1/2).\label{eq:riccati}
\end{equation}
\end{proposition}
\begin{proof}
Write $U=F_t'/F_t=u-iS$ on the line of height $p$.
Equation~\eqref{eq:poisson} and its differentiated sums give
\[ |u_x|\le S/d,\qquad |S_{xx}|\le2S/d^2. \]
The heat equation implies $U_t=-U_{xx}-2UU_x$, hence
\[ S_t+2uS_x=-S_{xx}-2Su_x. \]
Along $x'=2u$, this yields $S'\ge-2S/d^2-2S^2/d$.
The Riccati solution with initial value $C$ is the right side of
\eqref{eq:riccati}. The characteristics are complete in both time
directions: growth and a uniformly zero-free anchor $ip$ imply
$|u(x,p,t)|=O((1+|x|)\log(2+|x|))$ on bounded time intervals, which
prevents finite-time escape. Their flow is onto $\mathbb R$, proving
uniformity in $x$.
\end{proof}

At $p=3$, the right side is $m_C(t)$ in \eqref{eq:dm}. This is a
quantified loss, rather than an assertion that the time-zero lower bound
persists unchanged.

\section{The comparison when contacts are compact}
We briefly retain the confinement hypothesis used in Note~2: on each
$[\delta,T]$, $\delta>0$, the nonreal zeros lie in a bounded horizontal
region. This hypothesis will be established for approximants only.

At a highest simple zero $x+ih$, put $Q=h^2$. The conjugate-pair kernel
\[ K_h(d,v)=\frac{h-v}{d^2+(h-v)^2}
                 +\frac{h+v}{d^2+(h+v)^2} \]
obeys $K_h/h\ge K_p/p$ whenever $|v|\le h$ and $p^2\ge5h^2$.
Clearing denominators reduces this to nonnegativity of
\[ D^2+(A+B-6r)D+(A-r)(B-r),\quad
D=d^2,\ A=h^2,\ B=p^2,\ 0\le r=v^2\le A. \]
The zero-motion identity, with the zero's own conjugate pair removed,
therefore gives
\begin{equation}
 Q'\le-2-\frac{4Q}{p}S(x,p,t)+\frac{8Q}{p^2-Q}.
 \label{eq:highest}
\end{equation}
This form of the comparison appears in Hakimoglu-Brown \cite{hb}; the
algebra and its non-even extension are recorded in Notes~1--2.

Classical strip contraction gives $Q\le1$. At $p=3$,
$8/(9-Q)\le1$, so Proposition~\ref{prop:floor} gives
\[ Q'\le-2-\left(\frac43m_C(t)-1\right)Q. \]
The corresponding solution from one is
\begin{equation}
 \overline Q(t)=\frac{1-2J_C(t)}{I_C(t)}.\label{eq:qbar}
\end{equation}
Until it reaches zero it decreases, with descent rate at least
$2-\overline Q\ge1$.

For completeness, use a strictly higher comparison curve, with initial
value $1+\varepsilon$ and a positive $\varepsilon$ added to its
differential equation. This supplies an initial time gap uniformly in
space. Confinement makes a later first contact compact. A simple contact
contradicts \eqref{eq:highest}; a multiple contact has a backward Hermite
split above the smooth curve. Letting $\varepsilon\downarrow0$ yields
$M(t)\le\overline Q(t)$ while the latter is positive. If it reaches zero
at $T_0$, strip contraction makes $F_{t+\overline Q(t)/2}$ real-rooted for
$t<T_0$. Heat continuity and Hurwitz then make $F_{T_0}$ real-rooted;
strip contraction preserves that property afterwards.
The extended contact proof is Note~2, Section~4 \cite{note2}.

\section{Removing confinement by a real sine tail}\label{sec:approx}
Fix $0<C_1<C$, where $C$ is the input in
Theorem~\ref{thm:general}, and write
\[F(z)=c e^{az}z^m\prod_{\rho\ne0}E_1(z/\rho),\qquad
 E_1(w)=(1-w)e^w,\quad c,a\in\mathbb R.\]
Let $P_R$ truncate this product at $|\rho|\le R$, retaining conjugate
pairs and multiplicities. Set
\begin{equation}
 \kappa=2(C_1+1),\quad h=\pi/\kappa,\quad M=\lfloor R/(4h)\rfloor,
 \qquad G_R(z)=P_R(z)\prod_{k>M}\left(1-\frac{z^2}{k^2h^2}\right).
 \label{eq:approx}
\end{equation}
The new zeros are real, and $G_R\to F$ locally uniformly.

\begin{lemma}\label{lem:approx}
For all sufficiently large $R$, $S_{G_R}(x,3,0)\ge C_1$ for every real
$x$. Each fixed approximant flow has uniformly confined nonreal zeros
on every bounded time interval. Moreover, $(G_R)_t\to F_t$ locally
uniformly, uniformly on bounded nonnegative time intervals.
\end{lemma}
\begin{proof}
We give the estimates that carry the argument; their extended proofs
are in Note~3, Sections~2--4 \cite{note3}.
For $|x|\le R/2$, the omitted true-zero Poisson tail is at most
\[16(p+1)\sum_{|\rho|>R}|\rho|^{-2}=o(1).\]
For $|x|\ge R/2$, the full real lattice has lower bound
$\kappa\tanh(\kappa p)$, and removing its central zeros costs at most
$8p/(hR)+16p/R^2$. At $p=3$ this leaves a bound exceeding $C_1$;
the strict input gap $C-C_1$ handles the central region.

For fixed $R$, the sine product rewrites $G_R$ as
$e^{b_Rz}r_R(z)\sin(\kappa z)$ with $b_R$ real and $r_R$ real rational,
its poles canceled by sine zeros. If $r_R(z)\sim c_Rz^{d_R}$ at infinity,
Gaussian imaginary translation gives, uniformly for $|\Ima z|\le1$
and bounded $t$,
\begin{equation}
 \frac{(G_R)_t(z)}{c_Rz^{d_R}e^{b_Rz+(\kappa^2-b_R^2)t}}
 =\sin\bigl(\kappa(z-2b_Rt)\bigr)+O_{R,T}(|\Rea z|^{-1}).
 \label{eq:sine}
\end{equation}
Along the Gaussian line the real part remains $\Rea z$, so no rational
pole is crossed. Rouch\'e disks about the distant real sine zeros contain
exactly one zero each; reality forces it to be real. The complement in
the strip is zero-free. The strip itself follows from the genus-one
operator theorem of Cardon \cite[Theorem~1.5]{cardon}, with quadratic
parameter $2t$ and no quadratic exponential in $G_R$.

Finally, for any fixed $1<\sigma<2$, canonical-product bounds give
$|P_R(z)|\le e^{A(1+|z|)^\sigma}$ uniformly in $R$. The added tail is
bounded by $\sinh(\kappa|z|)/(\kappa|z|)\le e^{\kappa|z|}$.
These provide a common sub-Gaussian majorant in
\[ F_t(z)=\frac1{\sqrt\pi}\int_{\mathbb R}e^{-v^2}
                   F(z+2i\sqrt t\,v)\,\dd v. \]
Dominated convergence proves the heat-limit assertion. The limit is
nonzero: real Gaussian translation inverts imaginary Gaussian translation
under this same majorant (Note~3, Section~4).
\end{proof}

Apply the compact comparison to $G_R$. If $J_C(T)>1/2$, choose $C_1<C$
close enough that $J_{C_1}(T)>1/2$. The lemma makes $(G_R)_T$ real-rooted.
Hurwitz's theorem in each open half-plane passes this to the nonzero
limit $F_T$. For equality, first use $T+\eta<1/2$ and then let
$\eta\downarrow0$; the case $T=1/2$ is classical strip contraction.
This proves Theorem~\ref{thm:general}. No bound uniform in $R$ on the
approximants' confinement regions is required.

\section{Arithmetic input and exact scalar arithmetic}
The initial Dirichlet logarithmic derivative obeys
\begin{align}
 S_\chi(x,p,0)&=\frac12\Rea\frac{\xi_\chi'}{\xi_\chi}
                    \left(\frac{1+p+ix}{2}\right),\label{eq:initialidentity}\\
 S_\chi(x,p,0)&\ge C_{q,\mathfrak a}(p):=
 \frac14\left[\log\frac q\pi+
 \psi\left(\frac{1+p+2\mathfrak a}{4}\right)\right]
 -\frac12\left(-\frac{\zeta'}\zeta\right)\left(\frac{1+p}{2}\right).
 \label{eq:c}
\end{align}
Indeed, the Hermitian functional equation gives
\eqref{eq:initialidentity}; the gamma logarithmic derivative is bounded
below by its value on the real axis, and the Euler product gives
$|L'/L(s,\chi)|\le-\zeta'/\zeta(\Rea s)$ for $\Rea s>1$.
At $p=3$ the even case is the smaller of the two parities.

The finite certificate in Note~2, Section~5 \cite{note2}, proves
\begin{equation}
 C_{q,\mathfrak a}(3)>\frac{1221}{200}=6.105>6
 \qquad(q\ge10^{12},\ \mathfrak a=0,1).\label{eq:c6}
\end{equation}
It uses $\log10>23/10$, $\log4<7/5$, $\pi<4$,
$\gamma<H_{1000}-\log1000<29/50$, and a prime-power sum through 1000
with the positive tail $(\log1000+1)/1000$ to prove
$(-\zeta'/\zeta)(2)<3/5$. All decisions use exact rational enclosures.

The integral in \eqref{eq:ij} has the elementary form
\begin{equation}
 J_C(t)=\frac{3(1+2C)}{14C^2}(D_C(t)^{7/3}-1)
       -\frac3{20C^2}(D_C(t)^{10/3}-1).\label{eq:exactj}
\end{equation}
For $C=6$, $T=27/100$, the same certificate encloses
\[0.502464325582617834\le J_6(T)\le0.502464325582617835.\]
Alternating exponential series and rational cube-root bisection make
the lower endpoint rigorous. Theorem~\ref{thm:general} and
\eqref{eq:c6} prove the explicit part of Theorem~\ref{thm:dirichlet}.
The certificate programme, parameters, hashes, and extended scalar
verification are referenced from Notes~2--3 rather than reproduced here.
Its original record was conditional on an analytic contract; Section~
\ref{sec:approx} supplies that contract without changing the arithmetic.

For $0\le t\le1/2$,
\[J_C(t)\ge e^{-1/2}
\frac{(1+Ce^{-1/4}t)^{7/3}-1}{(7/3)Ce^{-1/4}}.\]
Consequently
\begin{equation}
 B(C):=\frac{e^{1/4}}C
 \left[\left(1+\frac76e^{1/4}C\right)^{3/7}-1\right]
 \label{eq:bc}
\end{equation}
is a sufficient reality time whenever $B(C)\le1/2$ and
$C\le C_{q,\mathfrak a}(3)$. Since
$C_{q,\mathfrak a}(3)=\frac14\log q+O(1)$ uniformly in parity,
$B(C)=O((\log q)^{-4/7})$ for the natural input. This proves the
remaining assertions of Theorem~\ref{thm:dirichlet}.

\section{Scope, controls, and provenance}
The sine-tail limit bypasses confinement for the original Dirichlet
flow; it does not prove that separate assertion. A periodic real-entire
heat flow may have a positive initial logarithmic-derivative floor and
unboundedly many nonreal zeros at positive times. The example and its
scope are in Note~2, Section~7. This does not contradict the present
quantitative comparison.

The arithmetic gain enters through \eqref{eq:c}. Euler-factor collision
controls can violate its precise Euler-product allowance. Weighted-square,
reflected, and primitive-compatible controls need not have the character
completion or that allowance. The general comparison remains valid when
their actual input floor meets its assumptions; no character-specific
lower bound is assigned to them by analogy. This is smoothing at
positive time, not a proof of GRH at time zero.

A literature correction is necessary. Ki--Kim--Lee
\cite[Section~6, Example~1]{kkl} already extend their qualitative
strict bound below $1/2$ to real primitive Dirichlet characters. The
initial investigation's search missed that precedent. Their example
for a complex character uses a real self-dual combination, rather than
the individual non-even family in \eqref{eq:heat}. We make no priority
claim for the quantitative results of this draft.

\paragraph{Acknowledgment and review status.}
This manuscript was drafted for Edward Baker with substantial assistance
from OpenAI's \texttt{gpt-6.1-sol} in Codex, recorded reasoning effort
\texttt{ultra}. Same-model subagents assisted with derivation and
adversarial proof checks. The exact scalar programme was also generated
with LLM assistance. These checks do not constitute independent
specialist review. Notes~1--3 preserve the mathematical dependencies,
corrections, and calculation records.

\begin{thebibliography}{9}
\bibitem{dobner} A. Dobner,
\emph{A proof of Newman's conjecture for the extended Selberg class},
Acta Arithmetica \textbf{201} (2021), 29--62;
\href{https://arxiv.org/abs/2005.05142v2}{arXiv:2005.05142v2}.
\bibitem{cardon} D. A. Cardon,
\emph{Complex zero strip decreasing operators},
\href{https://arxiv.org/abs/1401.0237v2}{arXiv:1401.0237v2}, Theorem~1.5;
J. Math. Anal. Appl. (2015),
\href{https://doi.org/10.1016/j.jmaa.2015.01.026}{doi:10.1016/j.jmaa.2015.01.026}.
\bibitem{kkl} H. Ki, Y.-O. Kim, and J. Lee,
\emph{On the de Bruijn--Newman constant}, Advances in Mathematics
\textbf{222} (2009), 281--306,
\href{https://doi.org/10.1016/j.aim.2009.04.003}{doi:10.1016/j.aim.2009.04.003}.
\bibitem{hb} C. Hakimoglu-Brown,
\emph{A conditional bound $\Lambda\le0.13963$ for the de Bruijn--Newman constant},
\href{https://arxiv.org/abs/2610.12803v1}{arXiv:2610.12803v1} (2026), Section~4.
\bibitem{note1} Research Note 1, prepared for E. Baker with Claude Opus 5.5,
\emph{Normalization and transfer contract}, 17 October 2026,
\href{../notes/round-1/01_NORMALIZATION_AND_TRANSFER_CONTRACT_CLAUDE_20261017.md}{notes/round-1}.
\bibitem{note2} Research Note 2, prepared for E. Baker with GPT-6.1-sol,
\emph{An explicit positive-time floor and a comparison conditional on confinement},
17 October 2026,
\href{../notes/round-2/02_BURGERS_FLOOR_AND_CONDITIONAL_COMPARISON_CODEX_20261017.md}{notes/round-2}.
\bibitem{note3} Research Note 3, prepared for E. Baker with GPT-6.1-sol,
\emph{Removing confinement from the strip comparison by a real sine tail},
17 October 2026,
\href{../notes/round-3/03_SINE_TAIL_LIMIT_AND_COMPARISON_CODEX_20261017.md}{notes/round-3}.
\end{thebibliography}
\end{document}
